\section{Three ways to choose the next brick}
\label{sec:generators}

\subsection{Greedy search (the baseline and the teacher)}
The baseline builds layer by layer from the ground up. Within a layer it repeatedly scores every legal placement of
every brick type and takes the best. The score puts \emph{hard} cells first: a target cell with nothing beneath it can
only ever be covered by a brick that also overlaps the layer below, so the search covers those before they become
unreachable. After hard cells it prefers placements that cover more new target cells, then fewer filler cells, then
more support underneath; ties break on the number of distinct bricks below (to interlock), part cost and position.
Three filler policies (none, at most a quarter, at most half) are run and the best validated IoU is kept. The whole
search takes 0.07\,s per shape on average. Its output, ordered by the planner, is also the \emph{teacher} that the
learned policy imitates.

\subsection{The policy network}
The learned generator is a candidate-scoring network (Figure~\ref{fig:policy}). Its input is the grid as three
voxel channels (target, current assembly, remaining target), encoded by four 3D convolutions with 32 channels at
full resolution (the grid is small enough that no pooling is needed). For each legal placement the encoder's features
are averaged over the placement's footprint cells and concatenated with 13 hand-computed features (the brick type as
a one-hot code, rotation, position, and the counts of new target, filler and supporting cells); a small MLP turns that
into one logit. A separate head scores the \emph{stop} action from a global pool of the feature map. A softmax over
all candidates plus stop gives the policy. Because only enumerated legal placements receive logits, legality is
inherent at training and at inference. The policy has 97k parameters; RL adds a value head and two inputs described
below, for 101k in total.

\begin{figure}[t]
  \centering
  \resizebox{\linewidth}{!}{%
  \begin{tikzpicture}[
    box/.style={draw=muted, fill=boxc, rounded corners=2pt, align=center, font=\small, inner sep=4pt, minimum height=8mm},
    arr/.style={-{Stealth[length=2mm]}, muted, line width=0.7pt},
    lab/.style={font=\scriptsize, text=muted, align=center}]
    \node[box] (vox) at (0,0) {voxel grid, $16\times16\times8$\\3 channels (SFT) / 4 (RL)};
    \node[box] (enc) at (4.0,0) {4\,$\times$ Conv3d(32) + GELU\\full resolution};
    \node[box] (pool) at (8.0,0) {mean over the candidate's\\footprint cells};
    \node[box] (mlp) at (11.6,0) {candidate MLP\\$64\!\to\!64\!\to\!1$};
    \node[box] (soft) at (14.6,0) {softmax over\\$N$ candidates\\+ STOP};
    \node[box] (feat) at (9.0,-2.0) {13 candidate features (+2 connectivity in RL):\\part one-hot, rotation, $x$, $y$, layer, new-target /\\filler / support counts; \emph{joins main}, \emph{merges two}};
    \node[box] (glob) at (4.0,-2.0) {global mean + max\\of the feature map};
    \node[box] (val) at (0.3,-2.0) {value head\\(RL only)};
    \node[box] (stop) at (14.6,-2.0) {STOP head};
    \draw[arr] (vox) -- (enc);
    \draw[arr] (enc) -- (pool);
    \draw[arr] (pool) -- (mlp);
    \draw[arr] (feat) -- (mlp);
    \draw[arr] (mlp) -- (soft);
    \draw[arr] (enc) -- (glob);
    \draw[arr] (glob) -- (val);
    \draw[arr] (glob.south) -- ++(0,-1.1) -| (stop.south);
    \draw[arr] (stop) -- (soft);
  \end{tikzpicture}}
  \caption{The candidate-scoring policy. Every legal placement is scored from pooled voxel features and a few
  hand-computed counts; STOP is scored from a global pool. 97k parameters (101k with the RL value head and
  connectivity inputs).}
  \label{fig:policy}
\end{figure}

\subsection{Imitation (SFT)}
\paragraph{Data.} Targets come from two sources. Ten procedural families (boxes, stacks, capped towers, bridges,
polyominoes, ellipsoids, cylinders, stepped shapes, blobs, rows of arches), 200 instances each, grounded and
face-connected. And a pool of Thingi10K meshes with every benchmark object excluded, voxelized by the production
pipeline. The supervision for each target is the greedy search's validated build in planner order, ending with the
stop action. We trained on two versions of this pool. \textbf{v1} kept 150 real meshes (at most 50k faces, and only
those on which the teacher covered at least half the target), so training was 93\,\% procedural. \textbf{v2} kept
every eligible mesh, 640 real meshes up to 200k faces, and draws half of each training batch from them. Examples
are split 80/10/10 by object, with all symmetric copies and all files of one Thingiverse object in the same split.

\paragraph{Training.} Teacher forcing with cross-entropy over candidates plus stop; AdamW at $10^{-3}$ with a cosine
schedule, batch 32, 20{,}000 steps, 8-fold symmetry augmentation (rotations and reflections of the grid). Decisions
are streamed by a background producer because materialising every augmented decision would need about 40\,GB.
Checkpoints are selected by free-running decode IoU on a fixed validation subset, not by teacher-forced accuracy,
because the latter hides compounding errors. One run takes 25 to 28 minutes on the laptop's GPU.

\subsection{Reinforcement learning}
Imitation reproduces the teacher's per-step choices but cannot see what only shows up in the finished build: whether
the layers interlock into one component, or whether an early placement blocked a cell for good. RL scores the
finished build directly.

\paragraph{Environment.} An episode places bricks on one target with the same candidate set plus stop, under
temperature-1 sampling. The state tracks a union-find over stud connections so the \emph{main} (largest grounded)
component is known at every step, exactly as the trim will decide it. Two things are added to the observation: a
fourth voxel channel marking the main component, and two per-candidate features, \emph{joins the main component} and
\emph{merges two components}. The SFT weights are zero-padded into this layout, so the initial RL policy produces the
SFT policy's logits exactly. The terminal reward is computed after the shared finalize step:
\[
  r \;=\; \iou_{\text{after trim}} \;+\; 0.1\cdot[\text{stable under the model}],
\]
with no per-step shaping. Validity cannot fail under the candidate rules, so no penalty is needed.

\paragraph{Algorithm.} PPO~\citep{schulman2017ppo} with a value head on the same encoder; advantage is reward minus
value (terminal reward only), normalised per batch; clip 0.2, entropy bonus 0.01, a KL penalty of 0.05 towards the
frozen initial policy, Adam at $10^{-4}$, gradient clipping 1.0. Masked actions are removed from the entropy and KL
terms before any multiplication; leaving them in poisons the backward pass with NaNs. Each iteration samples 128
episodes in 8 forked CPU workers and runs 3 epochs of updates on the GPU; 250 iterations take 1.6 to 2.4 hours. The
best checkpoint is chosen by greedy-decode IoU on a fixed 64-example validation subset that mirrors the training mix
(half real meshes in v2).

\paragraph{Decode rule.} At inference the policy decodes greedily (argmax). The stop action can be used as learned,
or \emph{masked}, in which case decoding continues until no legal placement remains or the cap is reached, exactly as
the greedy search stops. Section~\ref{sec:benchmark} describes how this choice was made.
